\documentstyle[%


righttag%


,11pt%


,amssymb%


,russian%               remove in Eng.


,izvru%                 Set izven% in Eng.


% ,izvru-ks             for brief communications


]{amsart}





%       Math definitions


%dot


\newcommand{\spc}[1]{\bold #1}


\newcommand{\eqdef}{\overset{\mathrm{def}}{=}{}}


\newcommand{\kr}{\operatorname{ker}}


\newcommand{\tts}[1]{}





%\hoffset=-15mm%                  For Eng.


\setcounter{page}{66}


\god{2000}


\nomer{4~(455)} %                        Remove in Eng.


%\nomer{Vol.\,44, No.\,4}%               Set in Eng.


%\str{\pageref{begin}--\pageref{end}}%  Set in Eng.


\udk{517.988}





\author{\it С.А.\,Гусаренко}


% NB:   ^^ remove in Eng.


\title{О минимизации квадратичных функционалов\\


с ограничениями в виде равенств}


\thanks{Работа выполнена при финансовой поддержке Российского фонда


фундаментальных исследований (99-01-01278, 96-15-96195)


и Конкурсного центра по исследованиям в области


фундаментального естествознания, Санкт-Петербург.}





\date{}


%\pagestyle{myheadings}%         Set in Eng.


\pagestyle{plain} %              Remove in Eng.


\begin{document}


%\thispagestyle{plain}%          Set in Eng.


\maketitle


\label{begin}





Широкие классы вариационных задач, содержащих, например,


функционалы с отклоняющимся аргументом, как правило,


неразрешимы методами классического вариационного


исчисления. Попытки выйти за рамки этих методов были


предприняты, например, в \cite{D}.


Для решения таких задач был разработан


метод редукции вариационных задач к задаче


минимизации функционала в некотором гильбертовом пространстве


\cite{AR}--\cite{G2}.


Ниже предлагается общий подход к такой редукции, основанный на


непосредственном изучении экстремальной задачи.





Напомним, что


действующий в гильбертовом пространстве $\spc{D}$


линейный самосопряженный ограниченный оператор


$U:\spc{D}\to\spc{D}$ называется {\it положительно определенным}


на пространстве $\spc{D}$, если


$\langle Uy,y\rangle\ge0$ при всех $y\in\spc{D}$;


{\it строго положительно определенным}, если


$\langle Uy,y\rangle>0$ при всех $y\in\spc{D}$, $y\ne0$;


{\it сильно положительно определенным},


если $\langle Uy,y\rangle\ge


\gamma\|y\|^2$ при некотором\linebreak $\gamma>0$ и всех $y\in\spc{D}$.


Положительная определенность оператора $U$ эквивалентна


неотрицательности его спектра. Сильная положительная


определенность оператора $U$ означает, что при некотором $\gamma>0$


все точки спектра оператора $U$ не меньше $\gamma$,


при этом существует обратный оператор $U^{-1}$ (\cite{LS}, c.\,249).





Известные условия разрешимости задачи


минимизации квадратичного функционала


\begin{equation}


\frac12\langle x,Ux\rangle-\langle x,f\rangle\longrightarrow\inf,


\end{equation}


где $f\in\spc{D}$


(напр., \cite{AR}, \cite{G}), сформулируем в следующем виде.





\begin{thm}


Точка $x_0\in\spc{D}$ является решением задачи $(1)$ тогда и только


тогда, когда $Ux_0=f$ и оператор $U$ положительно


определен на пространстве $\spc{D}$.


Если оператор $U$ строго положительно определен


на пространстве $\spc{D}$, то


задача $(1)$ имеет не более одного решения.


Если оператор $U$ сильно положительно определен


на пространстве $\spc{D}$, то оператор $U^{-1}$ обратим, и


задача $(1)$ имеет единственное решение $x_0=U^{-1}f$.


\end{thm}





Рассмотрим задачу минимизации квадратичного функционала


с линейными ограничениями в виде равенств


\begin{equation}


\frac12\langle x,Ux\rangle-\langle x,f\rangle\longrightarrow\inf,\quad


\ell x=\alpha,


\end{equation}


где компоненты линейного вектор-функционала


$\ell:\spc{D}\to\spc{R}^k$ линейно независимы и $\alpha\in\spc{R}^k$.





Обозначим через $\spc{D}_\ell$ ядро вектоp-функционала $\ell$ и


представим пространство $\spc{D}$ в виде прямой суммы


$\spc{D}=\spc{D}_\ell\oplus\spc{D}^\ell$,


где $\spc{D}^\ell$ --- пространство, ортогональное к пространству


$\spc{D}_\ell$.


Пусть $P:\spc{D}\to\spc{D}_\ell$ --- ортогональный проектор на


пространство $\spc{D}_\ell$, определяемый равенством


$P=I-\ell^*(\ell\ell^*)^{-1}\ell$ ($I$ --- единичный оператор).


Тогда


$\ov P=I-P=\ell^*(\ell\ell^*)^{-1}\ell$ --- ортогональный проектор


на пространство $\spc{D}^\ell$.


Пространство $\spc{D}^\ell$ является ядром оператора $P$


и совпадает с множеством


значений оператора $\ell^*:\spc{R}^k\to\spc{D}$,


сопряженного к вектоp-функционалу $\ell$.


Отсюда, в частности, следует, что


это пpостpанство конечномеpно и его размерность равна $k$.





\begin{thm}


Точка $x_0\in\spc{D}$ является решением задачи $(2)$


тогда и только тогда, когда


$$


Ux_0-f\in\spc{D}^\ell,\quad


\ell x_0=\alpha,


$$


и оператор $U$ положительно определен на пространстве $\spc{D}_\ell$.


Если оператор $U$ строго положительно определен на пространстве


$\spc{D}_\ell$, то задача $(2)$ имеет не более одного решения.


Если оператор $U$


сильно положительно определен на пространстве $\spc{D}_\ell$,


то задача $(2)$ имеет единственное решение


$x_0=PU_\ell^{-1}Pf+(I-PU_\ell^{-1}PU)\ell^*(\ell\ell^*)^{-1}\alpha$,


где $U_\ell=PU:\spc{D}_\ell\to\spc{D}_\ell$ --- сужение


оператора $U$ на подпространство $\spc{D}_\ell$.


\end{thm}





В некоторых случаях установить положительную


определенность оператора $U$ на пространстве $\spc{D}_\ell$


можно с помощью того факта, что


для любого положительно определенного оператора


$A:\spc{D}\to\spc{D}$ однозначно определяется


такой положительно определенный оператор $\sqrt{A}:\spc{D}\to\spc{D}$,


что $(\sqrt{A})^2=A$ (\cite{LS}, c.\,246).





\begin{thm}


Пусть оператор $U$ представим в виде


$U=A^{-1}-K$, где самосопряженный оператор


$A$ обратим и сильно положительно определен.


Если все точки спектра оператора $K_1=\sqrt{A}PKP\sqrt{A}$ не превосходят


единицы, то оператор $PUP$ положительно определен. Если операторы


$PUP$ и $A^{-1}-PA^{-1}P$ положительно определены, то


все точки спектра оператора $K_1$ не превосходят


единицы.


\end{thm}





\begin{cor}


Если $\|PKPA\|\le1$,


то оператор $PUP$ положительно определен.


\end{cor}





\begin{cor}


Пусть $U=I-K$. Тогда оператор $PUP$ положительно определен


тогда и только тогда, когда


все точки спектра оператора $PKP$ не превосходят единицы.


\end{cor}





Включение


\begin{equation}


Ux-f\in\spc{D}^\ell


\end{equation}


эквивалентно уравнению $P(Ux-f)=0$. Это


уравнение удобнее записывать в другом виде. Пусть $\spc{B}$ ---


некоторое гильбертово пространство


и $W:\spc{B}\to\spc{D}$ --- оператор, множество значений которого


совпадает с ядром функционала $\ell$, т.\,е. с пространством $\spc{D}_\ell$.


Тогда $\kr W^*=\spc{D}^\ell$ и


включение (3) эквивалентно уравнению


\begin{equation}


W^*(Ux-f)=0.


\end{equation}


Уравнения вида (4) составляют множество ``уравнений Эйлера''


задачи (2).


При этом положительная определенность оператора $PUP:\spc{D}\to\spc{D}$


эквивалентна положительной определенности


оператора $W^*UW:\spc{B}\to\spc{B}$.





Обозначим через $\spc{L}[a,b]$ пространство


функций $x:[a,b]\to\spc{R}$, суммируемых на $[a,b]$,


с нормой $\|x\|=\int\limits_a^b|x(t)|\,dt$, через


$\spc{L}_2[a,b]$ обозначим гильбертово пространство


функций $x:[a,b]\to\spc{R}$, суммируемых с квадратом на $[a,b]$,


со скалярным произведением


$\langle x,y\rangle=


\int\limits_a^bx(t)y(t)\,dt$ и через $\spc{W}_2^{(n)}[a,b]$


обозначим пространство таких $n$ раз


дифференцируемых функций, что


$x^{(n)}\in\spc{L}_2[a,b]$, и co скалярным


произведением


$$


\langle x,y\rangle=\sum\limits_{i=0}^{n-1}x^{(i)}(a)y^{(i)}(a)+


\int\limits_a^bx^{(n)}(t)y^{(n)}(t)\,dt.


$$





Рассмотрим классическую вариационную задачу


\begin{equation}


\frac12\int_a^b(\dot x^{2}(t)-p(t)x^{2}(t)-2g(t)\dot x(t))\,dt


\longrightarrow\inf,\quad


x(a)=\alpha,\quad x(b)=\beta,


\end{equation}


где $x\in\spc{W}^{(1)}_2[a,b]$, $p\in\spc{L}[a,b]$,


$g\in\spc{L}_2[a,b]$, $\alpha$, $\beta\in\spc{R}$.


Задача (5) запишется в виде (2), если положить


$$


(Ux)(t)=x(t)-x(a)-\int_a^bp(s)x(s)\,ds-


\int_a^t\int_s^bp(\tau)x(\tau)\,d\tau ds,\quad


f(t)=\int_a^tg(s)\,ds.


$$


Здесь $\ell x=(x(a),x(b))$,


$\ell^*(\gamma_1,\gamma_2)=\gamma_1+(1+t-a)\gamma_2$,


проектор на пространство $\spc{D}_\ell$ опpеделен pавенством


$$


(Px)(t)=x(t)-x(a)-(x(b)-x(a))\frac{t-a}{b-a},


$$


пространство $\spc{D}^\ell$ двумерно и состоит из линейных функций вида


$c_1+c_2t$, где $c_1$, $c_2$ --- произвольные постоянные.





Отметим, что условия однозначной разрешимости и


неотрицательности функции Грина двухточечной краевой задачи


$({\cal L}x)(t)\eqdef\ddot x(t)+p(t)x(t)=z(t)$, $x(a)=x(b)=0$,


обеспечивающие справедливость соответствующей теоремы


о дифференциальном неравенстве, привлекали внимание многих математиков.


В работе \cite{AD}


была доказана эквивалентность следующих утверждений:


\begin{itemize}


\item[a)] существует такая функция $v\in\spc{W}_2^{(2)}[a,b]$, что на $[a,b]$


выполняются неравенства $v(t)\ge0$, $({\cal L}v)(t)\le0$ и


$v(a)+v(b)-\int\limits_a^b({\cal L}v)(t)\,dt>0$;


\item[б)] любое нетривиальное решение уравнения ${\cal L}x=0$ имеет на $[a,b]$


не более одного нуля;


\item[в)] функция Коши $C(t,s)$ уравнения ${\cal L}x=0$ положительна


при $a\le s<t\le b$;


\item[г)] краевая задача ${\cal L}x=z$, $x(a)=x(b)=0$


однозначно разрешима в пространстве $\spc{W}_2^{(2)}[a,b]$,


причем ее функция Грина $G(t,s)$ отрицательна при $t,s\in(a,b)$;


\item[д)] спектральный радиус оператора


$(Wx)(t)=-\int\limits_s^b(G^-(t,s)p(s)+x(s))\,ds$,


где $G^-$ --- функция Грина краевой


задачи $\ddot \xi(t)-p^-(t)\xi(t)=f(t)$,


$\xi(a)=\xi(b)=0$ и $p(t)=p^+(t)-p^-(t)$, $p^+(t)\ge0$, $p^-(t)\ge0$,


меньше единицы.


\end{itemize}





Оказывается, что эти условия эквивалентны требованию однозначной


разрешимости задачи (5).





\begin{thm} Утверждения {\em а)--д)} эквивалентны утвеpждениям


\begin{itemize}


\item[1)] вариационная задача $(5)$ однозначно разрешима при всех


$g\in\spc{L}_2[a,b]$, $\alpha$, $\beta\in\spc{R};$


\item[2)] все собственные значения оператора


$(Ky)(t)=-\int\limits_a^bG_0(t,s)p(s)y(s)\,ds$,


$K:\spc{C}[a,b]\to\spc{C}[a,b]$, где $G_0$ ---


функция Грина краевой задачи $\Ddot y=z$, $y(a)=y(b)=0$,


принадлежат промежутку $(-\infty,1);$


\item[3)] решение задачи Коши ${\cal L}x=0$, $x(a)=0$, $\Dot x(a)=1$


положительно


на $(a,b]$.


\end{itemize}


\end{thm}





Пусть пространство $\spc{D}$ изоморфно прямому произведению


пространств $\spc{B}\times\spc{R}^n$, изоморфизм


${\cal J}=\spc{D}\to\spc{B}\times\spc{R}^n$


задается равенством $z=\delta x$, $\beta=rx$, и


${\cal J}^{-1}(z,\beta)=\Lambda z+Y\beta$,


где $\Lambda:\spc{B}\to\spc{D}$,


$Y:\spc{D}\to\spc{R}^n$, $\delta:\spc{D}\to\spc{B}$,


$r:\spc{D}\to\spc{R}^n$.





Задача


\begin{equation}


\sum_{i=1}^m\langle T_ix,T^ix\rangle+


\langle T_0x,g\rangle\longrightarrow\inf,\quad


\ell x=\alpha,


\end{equation}


где $T_i$, $T^i:\spc{D}\to\spc{B}$, $i=0,\dotsc,m$, ---


линейные ограниченные операторы,


изучалась в \cite{AR}, \cite{G}.


Она запишется в виде (2), если положить


$U=\sum\limits_{i=1}^m(T_i^*T^i+T^{i*}T_i)$,


$f=-T_0^*g$.





Если оператор $W:\spc{B}\to\spc{D}$ действует на


ядро функционала $\ell$, то уравнение ${\cal L}x=\varphi$,


где ${\cal L}=W^*U=\sum\limits_{i=1}^m(Q_i^*T^i+Q^{i*}T_i)$,


$\varphi=W^*f=-Q_0g$,


$Q_i=WT_i:\spc{B}\to\spc{B}$,


$Q^i=WT^i:\spc{B}\to\spc{B}$,


назовем (следуя \cite{AR}) уравнением Эйлера задачи (6).


Положительная


определенность опеpатоpа $PUP$ эквивалентна


положительной опpеделенности опеpатоpа


$H=\sum\limits_{i=1}^m(Q_i^*Q^i+Q^{i*}Q_i)$.





Методы постpоения опеpатоpа $W$,


существенно зависящие от


соотношения между $n$ и $k$, предлагались в \cite{AR}, \cite{G}.





Отметим, что любую квадратичную задачу (6) можно записать в виде


\begin{equation}


\frac12\big(


\langle  \delta x, Q\delta x+Arx\rangle+


\langle A^*\delta x+qrx,rx\rangle\big) -


\langle  \delta x,\varphi\rangle-


\langle  rx,\psi\rangle\longrightarrow\inf,\quad


\Phi\delta x+\Psi rx=\alpha,


\end{equation}


где линейные самосопряженные операторы


$Q:\spc{B}\to\spc{Q}$, $q:\spc{R}^n\to\spc{R}^n$


ограничены, операторы $A:\spc{R}^n\to\spc{B}$,


$\Phi:\spc{B}\to\spc{R}^k$, $\Psi:\spc{R}^n\to\spc{R}^k$ линейные,


$\varphi\in\spc{B}$, $\psi\in\spc{R}^n$, $\alpha\in\spc{R}^k$.


Задачу (7) можно также записать в матричном виде


$$


\frac12\langle{\cal J}x,{\cal U}{\cal J}x\rangle+


\langle {\cal J}x, F\rangle\longrightarrow\inf,\quad


\langle L,{\cal J}x\rangle=\alpha,


$$


где


${\cal U}=\begin{pmatrix} Q&A\\A^*& q\end{pmatrix}$,


$F=\begin{pmatrix}\varphi\\ \psi\end{pmatrix}$,


$L=\begin{pmatrix}\Phi\\ \Psi\end{pmatrix}$.


Здесь $U=\delta^*Q\delta+r^*A\delta+\delta^*A^*r+r^*qr


={\cal J}^{*}{\cal U}{\cal J}$,


$f=\delta^*\varphi+r^*\psi={\cal J}^{*}F$, $\ell=\Phi\delta+\Psi r=L{\cal J}$.


Следовательно, точка $x_0\in\spc{D}$ является решением задачи (7), если


при некотором $c\in\spc{R}^k$


$$


{\cal L}x_0=\varphi+\Phi^*c,\quad


\ell_0x_0=\psi+\Psi^*c,\quad


\ell x_0=\alpha,


$$


где ${\cal L}=Q\delta+Ar$, $\ell_0=A^*\delta+qr$


и опеpатоp $PUP$ положительно определен.


В частности, если матрица $\Psi$ обратима (регулярный случай \cite{G}),


то, полагая


$W=\Lambda-Y\Psi^{-1}\Phi$, получим, что точка $x_0\in\spc{D}$


является решением задачи (7) тогда и только тогда, когда


$$


({\cal L}-\Phi\Psi^{*-1}\ell_0)x_0=\varphi-\Phi\Psi^{*-1}\psi,\quad


\ell x_0=\alpha,


$$


и оператор $H=Q-A^*\Psi^{-1}\Phi-\Phi^*\Psi^{*-1}A+


\Phi^*\Psi^{*-1}q\Psi^{-1}\Phi$ положительно определен.








\begin{thebibliography}{9}





\bibitem{D}


Драхлин М.Е., Макагонова М.А. {\em Функционально--дифференциальные


уравнения Эйлера} // Функц.-дифференц. уравнения. -- Пермь, 1987.


-- С.\,12--18.


\bibitem{AR}


Azbelev N.V., Rakhmatullina L.F. {\em


Theory of linear abstract functional differential


equations and applications}. -- Tbilisi, 1996. -- V.\,8.


-- 101~с.


\bibitem{G}


Груздев А.А. {\em О редукции экстремальных задач к линейным


уравнениям в гильбертовом пространстве} // Изв. вузов. Математика. --


1993. -- \No\,5. -- С.\,36--42.


\bibitem{GG}


Гpуздев А.А., Гусаpенко С.А. {\em


О редукции вариационных задач к экстремальным задачам


без ограничений} // Изв. вузов. Математика.


-- 1994. -- \No\,6. -- С.\,39--50.


\bibitem{G2}


Гусаpенко С.А. {\em


О вариационных задачах с линейными ограничениями} // Изв. вузов. Математика.


-- 1999. -- \No\,2. -- С.\,30--44.


\bibitem{LS}


Люстерник Л.А., Соболев В.И. {\em Краткий курс функционального


анализа.} -- М.: Высш. школа, 1982. -- 271 с.


\bibitem{AD}


Азбелев Н.В., Домошницкий А.И. {\em О дифференциальном


неравенстве Валле-Пуссена} // Дифференц. уравнения. -- 1986. -- \No 12.


-- C.\,2041--2045.





\end{thebibliography}





\vskip 2cm





\post{\it Пермский государственный университет}{\it Поступила}


\post{}{24.05.1999}


\label{end}


\end{document}





